\documentclass[preprint,12pt,authoryear]{elsarticle}

\usepackage{amsmath,amssymb,amsthm}
\usepackage{booktabs}
\usepackage{array}
\usepackage{tabularx}
\usepackage{graphicx}
\graphicspath{{./}{figures/}}
\usepackage{enumitem}
\usepackage[colorlinks=true,urlcolor=blue,linkcolor=blue,citecolor=blue]{hyperref}
\emergencystretch=3em
\setlength{\tolerance}{2000}

\newtheorem{proposition}{Proposition}
\newtheorem{corollary}{Corollary}
\newcommand{\pos}[1]{\left[#1\right]_+}
\newcommand{\dconv}{\Delta b_{\mathrm{conv}}}
\newcommand{\Ups}{\Upsilon_c}
\newcommand{\Eceil}{E_{\mathrm{ceil}}}
\newcommand{\Acmin}{A_c^{\min}}
\newcommand{\tconv}{\tau_{\mathrm{conv}}}

\journal{Ecological Modelling}

\begin{document}

\begin{frontmatter}

\title{How a rising biocapacity total can conceal a shrinking ecological base: a two-land model}

\author[1]{Amin Abaee\corref{cor1}}
\ead{amin_abaee@ut.ac.ir}
\cortext[cor1]{Corresponding author. ORCID: 0000-0002-0019-1842.}
\affiliation[1]{organization={Independent Researcher}, city={Tehran}, country={Iran}}

\begin{abstract}
\footnotesize
Biocapacity is usually reported as one total, although land types differ in yield and regeneration. A rise in the total can therefore conceal a decline in the land that supports future production. We use a two-land model with delays in regeneration and population response to study this possibility. When the supply available to people falls short of demand, the model converts ecological-capital land into fast-provisioning land. If fast land contributes more annual biocapacity per hectare than capital land, the transfer adds to the total even as capital-land area falls. The total alone cannot reveal that change in composition.

In a fixed-demand comparison—with capital quality at its maximum, no accumulated debt or yield boost, and restoration switched off—the reduced model has at most one balance point. The balance curve meets the minimum-area boundary at a baseline demand ceiling of about $1.14$\,gha without turning back (no fold), and the reduced system's slowest return rate stays nonzero (no critical slowing down). Above that ceiling, capital-land area approaches its minimum. When a shortfall disappears, conversion stops; without restoration, the lower area persists but does not keep declining. In a representative simulation, capital-land area falls $46.2\%$ and its modeled contribution to biocapacity falls $51.7\%$, while total model biocapacity rises $39\%$. These are conditional model results, not measurements of global physical stocks.

In the National Footprint and Biocapacity Accounts, Edition 2026 (1961--2023), total world biocapacity rises $24.9\%$, while biocapacity per person falls $52.7\%$ and the production footprint per person rises $16.7\%$. Cropland adds $24.06$ percentage points to the total change; Forest biocapacity subtracts $5.95$ points, and built-up land adds $7.44$ as an accounting class. No comparable global physical-stock series was assembled, so we report no physical-stock composite or decline finding. The SI specifies requirements for such a comparison but contains no physical-stock inventory. These results show why component values should accompany aggregate totals.
\end{abstract}

\begin{keyword}
\footnotesize
biocapacity \sep ecological footprint \sep land-use conversion \sep aggregation \sep identifiability \sep critical slowing down \sep delay differential equations \sep carrying capacity
\end{keyword}

\end{frontmatter}

\section{Introduction}\label{sec:intro}

Ecologists use \emph{carrying capacity} for the population that a resource base can sustain. It is difficult to estimate, and many models treat it as a fixed ceiling \citep{Cohen1995}. But when demand exceeds renewal, use can damage the resources that support future production and change that ceiling. This paper examines a related measurement question: can a total biocapacity number rise while part of the productive land base is shrinking?

The National Footprint and Biocapacity Accounts (NFA) report biocapacity as annual biological supply and the Ecological Footprint as annual demand on biologically productive area. Both are expressed in global hectares (gha). A gha is a hectare adjusted to the world's average biological productivity for a particular year; its annual basis is implicit \citep{GFNStandards2009,GFNGuidebook2019}. This adjustment combines area and productivity in one account value. A rise in the total cannot by itself show whether land became more productive, more land entered production, or the mix of land classes changed. Research on composite indicators calls this \emph{compensation}: a surplus in one component can offset a deficit in another \citep{Nardo2008,Becker2017,Fischer2022}. Related critiques address how the Ecological Footprint should be interpreted \citep{Blomqvist2013,GiampietroSaltelli2014,vdBerghGrazi2015,Galli2016}. Aggregate biocapacity remains a useful summary, but it should be read alongside its component values.

Human--environment models often summarize environmental change as a single stock or an overall balance. Examples include the World3 system \citep{Meadows1972,Meadows1974,Meadows2004}, models of Easter Island \citep{BranderTaylor1998,Anderies2000}, and coupled Earth-system frameworks \citep{Motesharrei2016,Donges2017,CalvinBondLamberty2018,HendersonLoreau2018}. Other studies debate whether human carrying capacity is fixed, variable, or a useful concept at all \citep{Cohen1995,vdBerghRietveld2004,Simon1996,Meyer1994,MeyerAusubel1999,Dasgupta2021}. Global demographic assessments often project population separately from environmental constraints \citep{UN2022}. Meanwhile, land-use studies document the conversion of forests, grasslands, and wetlands to cropland \citep{Haberl2007,Krausmann2013,Erb2017}. Together, this work establishes that land conversion matters. What remains unclear is what a single biocapacity total reveals when land classes differ in yield and regeneration.

Consider an orchard with two productive kinds of ground. Fast-provisioning ground—such as cropland and intensively managed systems—can yield a substantial annual harvest and be replenished from reserve land. Ecological-capital ground—such as forest, soil, wetland, and long-rotation systems—holds assets whose productive value and regenerative capacity are rebuilt more slowly. The harvest is an annual flow; the trees, soil, and other assets are stocks that support future harvests. When available annual provision does not meet demand, the model allows a shortfall to prompt conversion of ecological-capital ground into fast-provisioning ground. In orchard terms, meeting demand by cutting trees can raise the current harvest while drawing down the standing productive stock and may reduce later harvests. Because the aggregate harvest combines contributions from both productive books, this conversion can raise the total when fast ground contributes more per hectare, even as the capital book shrinks. Replanting can restore next year's crop without restoring the standing stock or all the conditions that support it. Neither orchard nor harvesters respond instantly: tree loss, soil recovery, and demographic adjustment unfold on different timescales. A stable or rising harvest can therefore coexist with a shrinking productive base. This analogy explains an information limit; it is not evidence that NFA accounts measure global physical stocks or a prediction of inevitable or abrupt ecosystem failure. In the two-book analogy, one record follows the annual harvest and the other follows the assets behind it.

We represent this distinction with two effective land classes. Fast-provisioning land—such as cropland and intensively managed systems—can provide a high annual yield. Ecological-capital land—such as forest, soil, wetland, and long-rotation systems—supports future production but regenerates differently and often more slowly. When a shortfall triggers conversion from capital land to fast land, total biocapacity can rise even as capital-land area $A_c$ falls. This is a conditional model result: NFA account categories do not map directly to $A_c$ or measure the physical stocks in the orchard analogy. Related material-ledger work calls the broader flow--stock masking a \emph{productivity illusion}; this paper studies its composition-specific form \citep{Abaee2026c}.

A second analogy concerns what a service measure can tell us. An elevator rated for ten people may keep running for a time while carrying fourteen, but continued service does not show whether repeated overloading has reduced the cable's safety margin. The analogy separates observed output from the condition of the asset behind it. It does not predict an inevitable or sudden failure: in the reduced fixed-demand model, the balance curve reaches a specified area minimum without a fold.

We ask three questions. First, when can conversion raise total biocapacity while capital-land area $A_c$ falls, how long can this pattern last, and can the total reveal the area loss? Second, does the modeled drawdown behave like a tipping point—with a fold, hysteresis, or early-warning signals—or does it follow a different path? Third, what limits the drawdown, and what would allow the land area to recover?

The paper makes three contributions. First, it builds a two-land model in which area is conserved, capital-land quality changes separately, and shortfalls can trigger land conversion with a delay. Second, it gives the exact condition under which conversion raises total biocapacity and shows why the total alone cannot reveal how land is divided between the two classes. It also separates a unique balance point under fixed demand from the history-dependent path the model takes to reach it. Third, it decomposes component changes in NFA Edition 2026 and examines a separate crop-area/yield proxy, while stating clearly what these accounts can and cannot measure. The paper also gives five testable predictions. Parameters are representative, not fitted, so simulated paths illustrate conditional mechanisms rather than forecasts.\section{Methods}\label{sec:model}

\subsection{A two-land model with conserved area and a separate quality state}\label{sec:equations}

We divide total bioproductive area $A_{tot}$ among three land classes. Fast-provisioning land $A_f$ represents cropland and intensively managed systems. Ecological-capital land $A_c$ represents forest, soil, wetland, and long-rotation systems as one effective class, rather than as a single real land type. Reserve land $A_r$ is fallow, retired, or not yet allocated. The areas always sum to $A_{tot}$. Capital land also has a quality state $q\in[0,q_{\max}]$ that affects how much annual biocapacity it provides; $q$ is not an area. Quality follows a logistic growth rule:
\begin{equation}\label{eq:Gc}
G_c(q)=\rho_c\,q\left(1-\frac{q}{q_{\max}}\right),
\end{equation}
We calculate the annual value of capital land per hectare as its direct yield plus a term that depends on quality growth:
\begin{equation}\label{eq:Ups}
\Ups(q)=b_c+b_{G,c}\,G_c(q),
\end{equation}
Here $b_c$ is the direct yield, and $b_{G,c}$ converts quality growth into added annual value. Total biocapacity is the sum of the fast-land and capital-land contributions
\begin{equation}\label{eq:B}
B=Y_f+Y_c=b_f A_f+\Ups(q)\,A_c ,
\end{equation}
The resulting account amount $B$ is in gha, with the annual accounting basis implicit \citep{GFNStandards2009,GFNGuidebook2019}. Demand is $E=eP$: population $P$ multiplied by the constant per-person footprint $e$ (gha per person; $\mathrm{cap}$ denotes one person in the equations and tables). The shortfall $S$ is the part of demand not met by the shares of the two land contributions available to people:
\begin{equation}\label{eq:S}
S=\pos{E-\sigma_f Y_f-\sigma_c Y_c},
\end{equation}
Here $\pos{x}=\max(x,0)$, and $\sigma_f=\sigma_c=1$ in the baseline. When a shortfall remains, the model converts capital land into fast-provisioning land at a finite rate:
\begin{equation}\label{eq:uc}
u_c=\frac{1}{\tconv}\min\!\left(\frac{S}{b_f},\;A_c-\Acmin\right)\qquad[\mathrm{ha\,yr^{-1}}],
\end{equation}
Here $S/b_f$ is the area needed to cover the shortfall at the fast-land yield, $\Acmin$ is the minimum capital-land area allowed, and $\tconv$ sets how quickly conversion occurs. Land moves between classes through one-to-one transfers:
\begin{align}
\dot A_c&=-u_c+R_{rc},\label{eq:Ac}\\
\dot A_f&=u_c+\mu A_r-\eta_f A_f,\label{eq:Af}\\
\dot A_r&=-\mu A_r+\eta_f A_f-R_{rc},\label{eq:Ar}
\end{align}
Here $\mu A_r$ moves reserve land into fast use (for example, through replanting), and $\eta_f A_f$ moves fast land back into reserve. The only route from reserve back to ecological-capital land is the surplus-gated restoration flow $R_{rc}$:
\begin{equation}\label{eq:Rrc}
R_{rc}=\chi_r\,A_r\,\pos{B-E},
\end{equation}
These transfer equations conserve total area, so reserve land is the remainder after $A_f$ and $A_c$ are set. Capital-land quality responds to its earlier value after a regeneration delay. The model also tracks accumulated shortfall as debt, which lowers fast-land yield. Population follows a delayed logistic law \citep{Pearl1925,Hutchinson1948}, with carrying capacity $K=\max(B/e,K_{\min})$:
\begin{align}
\dot q&=G_c\big(q(t-\tau_g)\big),\label{eq:q}\\
\dot P&=rP\left(1-\frac{P}{K(t-\tau_p)}\right),\qquad K=\max(B/e,\,K_{\min}),\label{eq:P}\\
b_f&=\big(b_{f0}+T_b(t)\big)\,e^{-\alpha D},\label{eq:bf}\\
\dot D&=\pos{E-B}-\eta D,\label{eq:D}\\
T_b(t)&=\omega_{Tb}\,\frac{\Delta b}{1+e^{-\kappa_w(t-t_{wave})}},\qquad \omega_{Tb}\in\{0,1\}.\label{eq:Tb}
\end{align}
The model keeps area, annual account amounts, and accumulated shortfall distinct. Land areas $A_f$, $A_c$, and $A_r$ are in hectares; $B$, $E$, and $S$ are annual account amounts in gha, with the annual basis implicit. $D$ accumulates the shortfall through time, so it is measured in gha\,yr, while $\dot D$ and the repayment term $\eta D$ are in gha. This debt is an internal model quantity, not an area or an item reported by the NFA.

The optional technology boost $T_b$ raises fast-land yield gradually, following a bounded logistic curve. It is off in the baseline, main figures, and tables; it is switched on only in the matched comparison in SI~\S S8. The model tracks fast area $A_f$, capital area $A_c$, capital quality $q$, population $P$, and accumulated shortfall $D$. Timber liquidation and direct restoration from fast land to capital land are possible extensions, but are off in all reported runs (SI~\S S1).

The conversion switch is exact: it turns off when the shortfall reaches zero, and land moves at a finite rate set by $\tconv$. Changes in $q$ alter value per hectare without moving land. Reserve replanting can maintain fast land when conversion is inactive; without replanting, retirement would drive fast-land area to zero at a no-conversion equilibrium.

The model has two reinforcing feedback loops. First, if a hectare of fast land contributes more annual biocapacity than a hectare of capital land ($b_f>\Ups(q)$), converting land raises $B$ at that moment. More biocapacity can support a larger population; after the demographic delay, higher population can raise demand and create another shortfall. Other land flows and quality changes also affect $B$, so this conversion effect alone does not determine whether total biocapacity rises or falls.

Second, a shortfall adds to debt $D$, while repayment reduces it. More debt lowers fast-land yield $b_f$, which can reduce $B$ and deepen the shortfall. With fixed demand and a continuing shortfall, debt approaches a finite level.

\subsection{Parameters and their status}\label{sec:params}

Table~\ref{tab:params} lists the baseline values and the evidence for each. The model groups forests, soils, wetlands, and other slow-renewing assets into one effective capital-land class, even though they differ in area, condition, yield, and recovery time. Most values are representative, not fitted. Published records show different recovery times across systems: forest biomass reaches about half of its reference in a median of $30$\,yr, and collapsed fish stocks in about $5.6$\,yr; soil organic carbon returns toward a new balance over roughly $17$--$23$\,yr (SI~\S S5.5) \citep{Poorter2016,HutchingsReynolds2004,Neubauer2013,Poeplau2011}. These observations provide timescale context; they do not establish a global average, estimate the model's regeneration rate $\rho_c$, or determine the delay $\tau_g$.

The NFA cropland account value divided by cropland area, and the separate cereal-yield/cropland-area proxy, describe account data. Neither is a fitted fast-land yield $b_f$ or a direct measure of ecological capital (SI~\S S5). Their boundaries do not establish a one-to-one match with model fast land. The baseline has $b_f-\Ups(q)>0$, so a transfer raises biocapacity at that moment; this is a conditional model setting, not an empirical prediction.

\begin{table}[!htbp]
\centering\footnotesize
\renewcommand{\arraystretch}{0.95}
\caption{Baseline parameter values and evidence status. Representative values are effective, lumped quantities; the NFA account/area quotient is descriptive, not a physical-yield measurement. Extension parameters are active only in scenarios specified in this table or the SI. Symbols and units are listed in SI~\S S1.}
\label{tab:params}
\setlength{\tabcolsep}{4pt}
\begin{tabularx}{\textwidth}{@{}>{\raggedright\arraybackslash}p{0.13\textwidth} >{\raggedright\arraybackslash}p{0.2\textwidth} >{\raggedright\arraybackslash}X >{\raggedright\arraybackslash}p{0.26\textwidth}@{}}
\toprule
Parameter & Value & Meaning & Status \\
\midrule
$\rho_c$ & $0.08$\,yr$^{-1}$ & quality regeneration rate of capital land & representative; not fitted \\
$b_{f0}$ & $0.85$\,gha\,ha$^{-1}$ & base fast-land yield per hectare & representative (2026 NFA account value/area quotient $\approx1.045$ in 1961; not a physical-yield measure) \\
$b_c$ & $0.05$\,gha\,ha$^{-1}$ & capital-land direct yield & representative \\
$b_{G,c}$ & $0.8$\,gha\,yr\,ha$^{-1}$ & regeneration-value coefficient for capital land & representative \\
$q_{\max}$ & $1.0$ (dimensionless) & normalized quality ceiling & normalization \\
$\mu$ & $0.06$\,yr$^{-1}$ & cropland maintenance (reserve$\to$fast) & representative \\
$\eta_f$ & $0.05$\,yr$^{-1}$ & fast-land retirement (fast$\to$reserve) & representative \\
$\tconv$ & $1.0$\,yr & conversion time constant & representative \\
$\alpha$ & $0.03$\,gha$^{-1}$\,yr$^{-1}$ & degradation of $b_f$ per unit debt & representative \\
$\eta$ & $0.05$\,yr$^{-1}$ & debt-repayment rate ($>0$) & representative \\
$e$ & $0.55$\,gha\,cap$^{-1}$ & per-capita footprint (constant) & representative \\
$r$ & $0.02$\,yr$^{-1}$ & intrinsic population growth & representative \\
$\tau_g,\ \tau_p$ & $20,\ 25$\,yr & regeneration and demographic delays & representative; $\tau_p$ lumps institutional lags \\
$A_{tot}$ & $2.5$\,ha & total bioproductive area & normalization \\
$\sigma_f,\ \sigma_c$ & $1,\ 1$ (dimensionless) & human-available shares of annual account amounts & baseline ($<1$: reservation extension) \\
$\chi_r$ & $0.10$\,gha$^{-1}$\,yr$^{-1}$ & surplus-gated restoration coefficient & on with endogenous demand; off under frozen demand \\
$\Acmin$, $A_r^{\min}$,\newline $K_{\min}$ & $0.05$\,ha, $0.10$\,ha,\newline $0.02$\,cap & typed floors & domain constraints \\
$\Delta b$, $\kappa_w$,\newline $t_{wave}$ & $1.0$\,gha\,ha$^{-1}$, $0.1$\,yr$^{-1}$,\newline $40$\,yr & technology wave & off in baseline; active only in the matched SI S8 experiment \\
\bottomrule
\end{tabularx}
\end{table}

\subsection{Data and empirical measures}\label{sec:data}

We use the official National Ecological Footprint and Biocapacity Accounts, Edition 2026, results and metadata version 1.0 \citep{Lo2026NFA}. The main analysis covers the complete world series from 1961 to 2023; estimated values for 2024--2025 are excluded. The 1961--2022 sensitivity uses the same edition and population series, so both windows share one data vintage.

From the workbook's \texttt{world\_data} sheet, we use \texttt{biocap\_total\_gha} for total biocapacity $B$, \texttt{efp\_total\_gha} for the Global Ecological Footprint of Production $E_{\mathrm{prod}}$, and \texttt{pop\_1000\_persons} for UN PopStat population. We convert population from thousands to persons and calculate per-person values from the unrounded totals, rather than from the workbook's rounded per-person columns. Ratios, indices, and logarithmic changes have no units.

The component series retain the workbook labels Cropland, Grazing land, Forest biocapacity, Fishing grounds, and Built-up land. These are NFA account values, not measurements of biomass, soil condition, wetland extent, fish biomass, or the model's effective area $A_c$. Built-up land is an NFA accounting class and is not ecological capital. Accordingly, the land-type movements are used as accounting-composition signals only. Variable definitions and endpoint calculations are summarized in SI~\S S5.

As a separate check, we combine OWID's FAOSTAT-derived world cropland-area series (arable land plus permanent crops) with cereal yield \citep{OWIDCereal2026,OWIDCropland2026}. Cereal yield is a limited indicator of crop productivity, not a measure of all-crop output or an NFA yield factor. We scale the proxy to match the NFA cropland share of total biocapacity in 1961, then calculate $X_t=\hat b_{f,t}A_{f,t}$ and the remainder $C_t=B_t-X_t$. This is an assumed account proxy, not a measured split of physical capital. The remainder can include differences in boundaries, categories, and calibration; it is not $\Upsilon(q)A_c$. The proxy and its calibration sensitivity are reported in SI~\S S5. We infer no physical-stock trend from it.

The aggregate-versus-per-capita comparison uses the exact identity $\Delta\ln B=\Delta\ln(B/P)+\Delta\ln P$. This extensive/intensive split is an accounting identity, not a causal decomposition of biocapacity production.

\subsection{Numerical solution}\label{sec:numerics}

We solve the delay equations by advancing in time with the method of steps and an explicit first-order Euler scheme. The demonstration run and threshold bisection use a time step $\Delta t=0.05$\,yr; drawdown sweeps use $0.1$\,yr, with selected thresholds checked at $0.05$\,yr. The positive-part threshold is exact. State limits keep $A_c\ge\Acmin$, $A_f\ge0$, $q$ within $[0,q_{\max}]$, $P,D\ge0$, and $K\ge K_{\min}$. Before each run, histories of $B$ and $q$ are held at their initial values over the longest active delay. In the controlled wave experiment, both runs share the same no-wave history; the wave changes $B$ at the start and later changes the delayed carrying-capacity input. Because floor-proximity values depend on the step size and chosen tolerance, we report them as bands rather than sharp thresholds. The reported trajectories use delayed quality, $G_c(q(t-\tau_g))$; SI~\S S3 compares this with the current-quality version, and SI~\S S8 reports the matched wave experiment and step-size checks.
\section{Results}\label{sec:results}

\subsection{When conversion raises the total as capital land declines}\label{sec:mask}

At a fixed quality state $q$, the annual value of capital land per hectare, $\Ups(q)$, has a concave quadratic shape. It is highest at $q^*=q_{\max}/2$; under the baseline parameters this maximum is $0.066$\,gha\,ha$^{-1}$. Of that amount, $0.016$ comes from quality growth and $0.05$ from direct yield, so regeneration adds a modest contribution in this case.

The immediate effect of conversion depends on the difference between the annual account contribution of a hectare of fast land and the effective value of a hectare of capital land. We call this the \emph{composition diagnostic}:
\begin{equation}\label{eq:dconv}
\dconv(q,D)=b_f(D)-\Ups(q).
\end{equation}
At baseline parameter values, $\dconv$ is $+0.787$, $+0.784$, and $+0.794$\,gha\,ha$^{-1}$ at $q=0.3$, $0.6$, and $0.9$, respectively. All three values lie in the composition-premium regime ($\dconv>0$), where conversion contributes positively to the aggregate at the margin. At $\dconv=0$, the contribution is first-order neutral (composition-neutrality); at $\dconv<0$, it is negative (composition-discount). These names classify the conversion term alone: maintenance, retirement, quality change, and other terms still affect the net change in $B$. At fixed $q$, capital-land contribution is linear in $A_c$, so it has no interior maximum as capital-land area changes.

\begin{proposition}[Conversion, aggregate change, and what the total can reveal]\label{prop:mask}
Let $T_b\equiv0$, $R_{rc}\equiv0$ and hold $D$ fixed, so that $b_f$ is constant. Along any solution of \eqref{eq:B}--\eqref{eq:q},
\begin{equation}\label{eq:identity}
\dot B=\dconv(q)\,u_c+b_f\big(\mu A_r-\eta_f A_f\big)+b_{G,c}\,G_c'(q)\,A_c\,\dot q .
\end{equation}
\begin{enumerate}[label=(\roman*)]
\item During conversion ($u_c>0$), capital-land area falls when restoration is off. The conversion term in the rate of change of $B$ is $\dconv(q)u_c$; it is positive exactly when $\dconv(q)>0$. It determines the sign of the net change only when maintenance and retirement balance ($\mu A_r=\eta_fA_f$) and quality is not changing ($\dot q=0$).
\item In the reduced comparison with fixed demand and fixed debt, baseline availability shares, balanced maintenance and retirement, and no quality change, a positive $\dconv$ makes the shortfall fall as conversion proceeds. The shortfall reaches zero if the balance is reached first. If the minimum-area constraint binds first, $A_c$ approaches $\Acmin$ and a shortfall can remain. This result does not apply unchanged when the other terms in the full model are active.
\item At fixed $b_f>0$ and $q$, the total $B=b_fA_f+\Ups(q)A_c$ gives one equation for the two land areas $A_f$ and $A_c$. Near an interior state, a small change in $A_c$ can be offset by $\delta A_f=-(\Ups/b_f)\delta A_c$ and $\delta A_r=-\delta A_f-\delta A_c$, preserving both $B$ and total area as long as land bounds are respected. Thus nearby land allocations can produce the same total, and $B$ alone does not identify $A_c$. For a finite observed change $\Delta B$, the offset is $\Delta A_f=(\Delta B-\Ups\Delta A_c)/b_f$; it is feasible only when all resulting areas remain within their bounds, so the range depends on the starting state. If $A_f$ and $b_f$ are observed independently, the residual $C=B-b_fA_f=\Ups(q)A_c$ identifies their product, not the area and value separately. An independent observation of $A_c$ would give $\Ups(q)=C/A_c$. A single value of $\Ups$ cannot separate $b_c$, $b_{G,c}$, $\rho_c$, and $q$; even with informative independent observations of $q$, it identifies at most $b_c$ and the product $b_{G,c}\rho_c$, not $b_{G,c}$ and $\rho_c$ separately.
\end{enumerate}
\end{proposition}

\begin{proof}
Differentiate \eqref{eq:B} with $b_f$ constant and substitute \eqref{eq:Ac}--\eqref{eq:Af} with $R_{rc}=0$: $\dot B=b_f(u_c+\mu A_r-\eta_fA_f)-\Ups u_c+b_{G,c}G_c'(q)A_c\dot q$, which is \eqref{eq:identity}. (i) Follows from \eqref{eq:identity} with $R_{rc}=0$; when the maintenance--retirement imbalance and quality-growth terms vanish, only the conversion term remains. (ii) Under the reduced-case conditions in the statement, $S=E-B$ and $\dot S=-\dot B=-\dconv u_c<0$ on $S>0$. Since $S$ is bounded below by zero and $A_c$ is bounded below by $\Acmin$, conversion either closes the shortfall or is limited by the floor. (iii) At fixed $b_f$ and $q$, the observation map is linear; an interior feasible point admits the stated small perturbations along a level set of $B$, while finite changes must also obey the area bounds. The residual $C$ identifies only $\Ups(q)A_c$ unless $A_c$ is independently observed; the parameter combination follows from the definition of $\Ups(q)$.
\end{proof}

Figure~\ref{fig:model}a shows the baseline run with frozen demand $E=0.825$\,gha and restoration switched off. It starts with $A_f=0.6$\,ha, $A_c=1.5$\,ha, $A_r=0.4$\,ha, and $q=0.9$. During conversion, total model biocapacity $B$ rises from $0.594$ to $0.825$\,gha ($39\%$), capital-land area falls from $1.500$ to $0.808$\,ha ($46\%$), and fast-land area rises from $0.600$ to $0.923$\,ha.

The quality state $q$ rises from $0.9$ to $1.0$. Because it starts above $q^*=0.5$, the peak of $\Ups(q)$, its per-hectare value falls from $0.05576$ to $0.05000$\,gha\,ha$^{-1}$. The modeled contribution $\Ups(q)A_c$ therefore falls from $0.08364$ to $0.04038$\,gha ($51.7\%$). This decline is larger than the $46.2\%$ decline in area because per-hectare value also falls. The land transfers conserve area: $\Delta A_f+\Delta A_c+\Delta A_r=0$.

The change is fastest at first, then slows. In the deficit-side model, the approach to balance is exponential, with a slow rate $|\lambda_1|=0.054$\,yr$^{-1}$—an e-folding time of about $19$\,yr. Within $25$\,yr, $98\%$ of the rise in $B$ and $85\%$ of the decline in $A_c$ have occurred ($B=0.820$\,gha and $A_c=0.913$\,ha); $97\%$ of the area decline has occurred by year $50$. The simultaneous pattern $\dot B>0$ and $\dot A_c<0$ lasts about $149$\,yr at the numerical resolution used.

The long tail reflects reserve adjustment. As fast land retires to the reserve, its contribution falls and a small shortfall reopens, allowing more conversion until $A_r/A_f=\eta_f/\mu$. The duration is therefore not a fixed horizon: its scale depends on the slow rate $\lambda_1$ and its size on $\dconv$ and the initial shortfall (SI~\S S2.2).

This result limits what an aggregate can tell us. Total $B$ alone cannot show whether capital land is shrinking while another component grows. Independent measures of fast-land area $A_f$ and yield $b_f$ let us calculate the remainder $C=B-b_fA_f$, but that remainder identifies only the capital-land contribution—not its area and per-hectare value separately. We call a trajectory with rising $B$ and falling $A_c$ a \emph{composition mask}. A distinct, though potentially overlapping, pattern is the \emph{liquidation mask}: demand $E$ remains steady while $A_c$ falls. The composition mask describes an aggregate-versus-area signature; the liquidation mask describes a demand-versus-area signature. The composition diagnostic $\dconv$ tells us the sign of conversion's marginal contribution, not necessarily the net change in $B$ when other terms are active. A liquidation mask does not require $\dconv>0$. Observing $B$, $E$, and $A_c$ distinguishes the two defining signatures. Here, “liquidation mask” names this trajectory; it is separate from the optional timber-liquidation extension, which is off in all reported runs.

\begin{figure}[!htbp]
\centering
\includegraphics[width=\textwidth]{fig1_model_dynamics}
\caption{Model results at the baseline parameter set (Table~\ref{tab:params}). Account amounts $B$ and $E$ are in gha; land areas are in ha; population is in cap. (a) Frozen demand $E=0.825$\,gha, restoration off, $A_{f0}=0.6$\,ha, $A_{c0}=1.5$\,ha, and $q_0=0.9$. $B$ rises from $0.594$ to $0.825$\,gha as $A_c$ falls from $1.500$ to $0.808$\,ha and $A_f$ rises from $0.600$ to $0.923$\,ha. Shading marks the interval with $\dot B>0$ and $\dot A_c<0$ (about $149$\,yr; only the first $60$\,yr are shown; within $25$\,yr, $98\%$ of the rise in $B$ and $85\%$ of the fall in $A_c$ occur). (b) Fixed-demand balance relation $B_{eq}(A_c)$ and its calculated areas $A_c^*(E)=1.110$, $0.626$, $0.385$, $0.143$, and $0.070$\,ha at $E=0.70$, $0.90$, $1.00$, $1.10$, and $1.13$\,gha. The curve meets the minimum $A_c=0.05$\,ha at $\Eceil\approx1.138$\,gha; above this demand, no balance point is admissible. (c) Endogenous-demand runs with $\tau_g=20$\,yr, $\tau_p=25$\,yr, and $\chi_r=0.10$\,gha$^{-1}$\,yr$^{-1}$. For each starting area $A_{c0}$, the minimum area reached varies continuously and monotonically with initial population $P_0$ (initial demand $E_0=eP_0$). The floor is approached asymptotically; triangles mark the proximity values in Table~\ref{tab:model}, defined by an approach to within $10^{-6}$\,ha. No discontinuity is observed.}
\label{fig:model}
\end{figure}

\subsection{Fixed demand leads to a minimum-area boundary, not a fold}\label{sec:ceiling}

When population has adjusted to the model's carrying capacity, $P=K=B/e$, so demand $E=eP$ equals biocapacity $B$ and the shortfall is zero. Active conversion is therefore transient in the endogenous-population runs. To study how the balance changes with demand, we also hold $E$ fixed and analyze the reduced land system. The model imposes a minimum capital-land area $\Acmin$; this is a \emph{typed floor}, a limit on one specific land class rather than on the aggregate.

\begin{proposition}[Balance point, demand ceiling and affine deficit-side dynamics]\label{prop:ceiling}
For this reduced comparison, hold demand at $E$, switch off the technology boost and restoration, set debt $D=0$, and fix quality at $q=q_{\max}$. This quality is the attractor of \eqref{eq:q} for $q(0)>0$. Let $k=\mu/(\mu+\eta_f)$ and
\begin{equation}\label{eq:Beq}
B_{eq}(A_c)=b_{f0}\,k\,(A_{tot}-A_c)+b_cA_c=b_{f0}kA_{tot}-(b_{f0}k-b_c)A_c .
\end{equation}
\begin{enumerate}[label=(\roman*)]
\item With no conversion ($S=0$), reserve-to-fast maintenance balances fast-land retirement. This gives $A_f=k(A_{tot}-A_c)$ and $A_r=(1-k)(A_{tot}-A_c)$, where $k=\mu/(\mu+\eta_f)$. The resulting total is $B_{eq}(A_c)$. If $b_{f0}k>b_c$, this total decreases as $A_c$ increases, so there is at most one balance point $B_{eq}(A_c^*)=E$:
\begin{equation}\label{eq:Acstar}
A_c^*(E)=\frac{b_{f0}kA_{tot}-E}{b_{f0}k-b_c}.
\end{equation}
The algebraic solution is feasible when $B_{eq}(A_{tot})\le E\le\Eceil$, where $\Eceil=B_{eq}(\Acmin)$. It lies above the minimum area exactly when $E<\Eceil$. At lower demand, the algebraic root lies above $A_{tot}$, and every feasible no-conversion allocation already supplies at least $E$.
\item While shortfall is positive and conversion is not yet limited by the minimum-area constraint, the land subsystem is affine in $z=(A_f,A_c)$: $\dot z=Jz+c$, with
\begin{equation}\label{eq:J}
J=\begin{pmatrix}-\tconv^{-1}-\mu-\eta_f & -\mu-\dfrac{b_c}{b_f\tconv}\\[8pt] \tconv^{-1} & \dfrac{b_c}{b_f\tconv}\end{pmatrix},
\qquad \det J=\frac{\mu+\eta_f}{\tconv}\Big(k-\frac{b_c}{b_f}\Big),
\end{equation}
and $\operatorname{tr}J<0$. The matrix does not depend on demand or on the current state. Its determinant is positive exactly when $B_{eq}$ decreases, and its fixed point is the balance point in (i). For $B_{eq}(A_{tot})\le E<\Eceil$, trajectories that remain in this regime approach that point at rates set by the eigenvalues of $J$. For $E>\Eceil$, the balance point would lie below the allowed minimum: $A_c$ approaches $\Acmin$ asymptotically and the shortfall tends to $S^*=E-\Eceil>0$. Thus the reduced balance curve meets the minimum-area boundary; it does not turn back or form a fold. The leading eigenvalue does not approach zero at $E=\Eceil$, and the reduced equilibrium relation has no hysteresis. These conclusions hold under the stated assumptions and do not imply that the realized land allocation is independent of its history.
\item Above the demand ceiling, with degradation active, debt approaches a unique value $D^*(E)$ satisfying $\eta D=E-B_{eq}(\Acmin;D)$, where $B_{eq}(\cdot;D)$ uses yield $b_{f0}e^{-\alpha D}$. This result requires $\eta>\alpha b_{f0}k(A_{tot}-\Acmin)$; the equilibrium debt increases with demand.
\end{enumerate}
\end{proposition}

\begin{proof}
(i) With $S=0$ and $R_{rc}=0$, \eqref{eq:Af} gives $\mu A_r=\eta_fA_f$, hence $A_r/A_f=\eta_f/\mu$ and, with $A_r=A_{tot}-A_f-A_c$, $A_f=k(A_{tot}-A_c)$; substituting into \eqref{eq:B} with $q=q_{\max}$ ($\Ups=b_c$) and $D=0$ gives $B=B_{eq}(A_c)$, and $S=0$ requires $B\ge E$. Strict monotonicity and \eqref{eq:Acstar} follow from $b_{f0}k>b_c$. (ii) On the deficit side $u_c=(E-b_fA_f-b_cA_c)/(b_f\tconv)$ is affine in $z$; substituting into \eqref{eq:Ac}--\eqref{eq:Af} with $A_r=A_{tot}-A_f-A_c$ gives $J$ and a constant vector $c$ by direct differentiation; $\det J$ follows by expansion, and $\operatorname{tr}J=-\tconv^{-1}(1-b_c/b_f)-\mu-\eta_f<0$ for $b_c<b_f$. The fixed point solves $u_c=0$ and $\mu A_r=\eta_fA_f$, i.e.\ it is the balance point. Because $J$ is constant, convergence rates are constant along the branch. (iii) The map $D\mapsto\eta D+B_{eq}(\Acmin;D)-E$ is continuous, negative at $D=0$ when $E>\Eceil$, tends to $+\infty$, and is strictly increasing when $\eta>\alpha b_{f0}k(A_{tot}-\Acmin)\ge-\partial_DB_{eq}$; hence it has a unique root, which increases with $E$.
\end{proof}

The uniqueness result is about the reduced equilibrium relation, not path independence of realized states. In the baseline trajectory, $R_{rc}=0$ and the conversion flow $u_c$ has no reverse counterpart. When the shortfall closes, $A_c$ stops declining and approaches its lower balance value; it remains below its pre-overshoot level because this run contains no restoration flow. The quality state $q$ evolves separately, so $\Ups(q)$ can continue changing as $A_c$ approaches that value. Thus the model has a unique frozen-demand balance point under Proposition~\ref{prop:ceiling}'s restrictions, but a history-dependent post-overshoot land allocation; this is not perpetual decline at equilibrium and is not a general claim about real ecosystems.

At baseline, $k=0.5455$ and the slope $b_{f0}k-b_c=0.414$\,gha\,ha$^{-1}$. The demand ceiling is $\Eceil=1.138$\,gha. The predicted balance areas—$1.110$, $0.626$, $0.385$, $0.143$, and $0.070$\,ha at demands of $0.70$, $0.90$, $1.00$, $1.10$, and $1.13$\,gha—match the integrated equilibria (Fig.~\ref{fig:model}b). The eigenvalues of $J$ are real and negative: $\lambda_1=-0.0537$ and $\lambda_2=-0.998$\,yr$^{-1}$. Across the tested fast-land yields, $\lambda_1$ changes only from $-0.0522$ at $b_f=0.60$ to $-0.0556$\,yr$^{-1}$ at $b_f=2.00$ and never changes sign (Table~\ref{tab:model}a).

Above the ceiling, the calculated debt equilibria are $D^*=1.19$ and $5.38$\,gha\,yr at $E=1.158$ and $1.238$\,gha; numerical integrations give $1.21$ and $5.39$\,gha\,yr. The condition for a unique debt equilibrium holds with margin: $\alpha b_{f0}k(A_{tot}-\Acmin)=0.034$\,yr$^{-1}$, below $\eta=0.05$\,yr$^{-1}$. Thus debt saturates rather than growing without bound under fixed demand.

We use \emph{critical slowing down} in its standard sense: as a system approaches a fold, its slowest return rate tends to zero and the return time $T_{ret}\approx-1/\lambda_1$ grows. Here the linearization is constant along the reduced balance branch; $\lambda_1$ does not approach zero, and the return time stays near $19$\,yr as $E$ approaches $\Eceil$. The model's signal is instead $A_c$ approaching its minimum, which an aggregate-only observation cannot reveal.

With population responding to supply, the minimum capital-land area reached during a run varies continuously and monotonically with initial population $P_0$ for each starting area $A_{c0}$ (Fig.~\ref{fig:model}c). The floor is approached asymptotically, not crossed, because $u_c\le(A_c-\Acmin)/\tconv$. A larger starting area supports a larger initial population before the trajectory approaches the floor.

Table~\ref{tab:model}b reports the initial populations at which $A_c$ comes within $10^{-2}$, $10^{-4}$, or $10^{-6}$\,ha of its minimum. These are tolerance-based proximity values, not abrupt thresholds. At the tightest tolerance, the values are $P_0\approx2.41$, $2.70$, $2.98$, and $3.14$\,cap for starting areas $A_{c0}=0.40$, $0.70$, $1.00$, and $1.18$\,ha, respectively. They rise with starting area, and no jump appears at any of the tested tolerances.

For $A_{c0}=1.0$\,ha, the $10^{-6}$\,ha proximity value changes by at most $0.001$\,cap across regeneration delays $\tau_g=5$, $20$, and $40$\,yr (SI Table S2). This narrow comparison does not establish that delay has no effect outside the tested range. The threshold is also the same to three decimal places with surplus-gated restoration on or off (SI~\S S4), because that restoration is inactive while $B<E$.

\begin{table}[!htbp]
\centering\small
\caption{Model results at the baseline parameter set. (a) Composition diagnostic at $q^*=q_{\max}/2$ ($\Ups(q^*)=0.066$\,gha\,ha$^{-1}$), eigenvalues of the deficit-side matrix $J$ of \eqref{eq:J} and the demand ceiling, for a sweep of the fast-land yield $b_f$ (all other parameters at baseline). (b) Initial population $P_0$ (cap; $E_0=eP_0$ in gha, with $e=0.55$\,gha\,cap$^{-1}$) at which the minimum capital land reached during an endogenous-demand run ($A_{f0}=0.6$\,ha, dimensionless $q_0=0.9$, $\tau_g=20$\,yr, $\tau_p=25$\,yr) comes within the stated distance of the typed floor $\Acmin=0.05$; the first two columns are read from a $P_0$ grid of step $0.1$ ($\Delta t=0.1$\,yr, $T=400$\,yr), the last is a bisection using the $10^{-6}$\,ha floor-proximity criterion ($\Delta t=0.05$\,yr, $T=1000$\,yr). Drawdown depth is continuous in $P_0$ (Fig.~\ref{fig:model}c); the thresholds are a band, not a single value.}
\label{tab:model}
\begin{tabular}{l l l l l}
\toprule
\multicolumn{5}{l}{(a) Fast-land yield sweep (all rows composition-premium, $\dconv>0$)}\\
$b_f$ (gha\,ha$^{-1}$) & $\dconv(q^*)$ (gha\,ha$^{-1}$) & $\lambda_1$ (yr$^{-1}$) & $\lambda_2$ (yr$^{-1}$) & $\Eceil$ (gha) \\
\midrule
$0.60$ & $+0.534$ & $-0.0522$ & $-0.975$ & $0.804$ \\
$0.85$ (baseline) & $+0.784$ & $-0.0537$ & $-0.998$ & $1.138$ \\
$1.20$ & $+1.134$ & $-0.0547$ & $-1.014$ & $1.606$ \\
$2.00$ & $+1.934$ & $-0.0556$ & $-1.029$ & $2.675$ \\
\bottomrule
\end{tabular}

\medskip
\begin{tabular}{l l l l}
\toprule
\multicolumn{4}{l}{(b) Initial population $P_0$ (cap) at which $\min_tA_c$ approaches $\Acmin$ to within:}\\ 
$A_{c0}$ (ha) & within $10^{-2}$\,ha & within $10^{-4}$\,ha & within $10^{-6}$\,ha \\
\midrule
$0.40$ & $\approx1.8$ & $\approx2.1$ & $\approx2.41$ \\
$0.70$ & $\approx2.2$ & $\approx2.4$ & $\approx2.70$ \\
$1.00$ & $\approx2.6$ & $\approx2.8$ & $\approx2.98$ \\
$1.18$ & $\approx2.8$ & $\approx3.0$ & $\approx3.14$ \\
\bottomrule
\end{tabular}
\end{table}

\subsection{Technology gains are finite; debt depends on the shortfall}\label{sec:bounded}

The technology boost and accumulated debt both affect fast-land yield $b_f$, but in different ways. The boost has a fixed upper limit. Debt, by contrast, can continue to rise while a shortfall persists, although it approaches a finite level when the shortfall itself settles.

\begin{proposition}[The technology boost is capped; debt depends on the shortfall]\label{prop:bounded}
Under \eqref{eq:bf}--\eqref{eq:Tb}:
\begin{enumerate}[label=(\roman*)]
\item The technology factor $(b_{f0}+T_b(t))/b_{f0}$ stays between $1$ and $1+\Delta b/b_{f0}$. Its log-scale gain is therefore bounded by $\ln(1+\Delta b/b_{f0})$ for all time ($0.78$ at baseline).
\item Debt satisfies $D(t)\le\max\{D(0),\ \sup_{s\le t}\pos{E-B}(s)/\eta\}$. If the shortfall settles at $S^*$, then $D$ tends to $S^*/\eta$. The resulting reduction in fast-land yield depends on the size of that sustained shortfall.
\item The debt-related yield loss exceeds the largest possible technology gain once the sustained shortfall exceeds $S^\dagger=(\eta/\alpha)\ln(1+\Delta b/b_{f0})$, which is $1.30$\,gha at baseline.
\end{enumerate}
With additive degradation, $b_f=b_{f0}+T_b-\alpha D$, the two effects enter symmetrically; the comparison in (iii) then reduces to $\Delta b$ versus $\alpha D^*$. The asymmetry above follows from the model's multiplicative form.
\end{proposition}

\begin{proof}
(i) is $0\le T_b\le\Delta b$. (ii) From \eqref{eq:D}, $\dot D\le\bar S-\eta D$ with $\bar S=\sup_{s\le t}\pos{E-B}$, so $D(t)\le\max\{D(0),\bar S/\eta\}$ by comparison; if $\pos{E-B}\to S^*$ then $D\to S^*/\eta$. (iii) compares $\alpha S^*/\eta$ with $\ln(1+\Delta b/b_{f0})$.
\end{proof}

With fixed demand and a persistent shortfall, $D$ approaches $S^*/\eta$; a larger shortfall gives a higher debt level. At baseline, the crossover shortfall $S^\dagger\approx1.3$\,gha exceeds baseline biocapacity $B$. In this parameter set, the balance with the minimum-area constraint limits conversion before debt-related yield loss exceeds the largest possible technology gain, although debt still lowers fast-land yield over time.

A technology boost can raise $B$ and carrying capacity; after the population delay, demand may rise as well. Its effect on debt depends on starting conditions and on which part of the run is summarized. In the matched experiment in SI~\S S8, for $P_0=1.0$\,cap, the largest debt over $0$--$250$\,yr is $0.02020$\,gha\,yr with the boost and $0.00813$\,gha\,yr without it (ratio $2.49$). For $P_0=1.2$\,cap, the largest value over the full run is lower with the boost ($0.05104$ versus $0.06122$\,gha\,yr), but the largest value after the wave midpoint at $40$\,yr is higher ($0.02518$ versus $0.01085$\,gha\,yr; ratio $2.32$). The effect therefore changes with initial population and with the time window used. These are conditional model results, not forecasts or empirical estimates.

\subsection{Recovery requires an explicit, gated restoration flow}\label{sec:recovery}

Conversion moves land from capital use to fast use; the model has no return path unless the restoration flow $R_{rc}$ is switched on. Fast-land retirement sends land to reserve, not back to capital land. Thus, reducing population, footprint, or shortfall does not by itself reverse earlier conversion. When $\chi_r=0$, capital-land area remains at its post-conversion level.

In one endogenous-demand run ($A_{c0}=1.0$\,ha, $P_0=2.0$\,cap), $A_c$ falls to $0.398$\,ha within four years and stays there without restoration. With surplus-gated restoration ($\chi_r=0.10$\,gha$^{-1}$\,yr$^{-1}$), it reaches $0.828$\,ha; the reported recovery fraction is $\mathcal R_c=(A_c(T)-A_c^{\mathrm{deg}})/(A_c^{\mathrm{init}}-A_c^{\mathrm{deg}})=0.71$ (SI~\S S4.3).

The two restoration gates act at different times. A surplus-gated flow is inactive during shortfall, but can continue after a surplus opens while population adjusts; reaching a specified target therefore requires a target-limited rule (SI~\S S4.4). A deficit-gated flow acts during shortfall but removes productive area when supply is already insufficient. At the tested rates, conversion exceeds restoration and the drawdown continues. These outcomes depend on timing and rates; neither gate is universally preferable.

\subsection{NFA account composition, 1961--2023}\label{sec:empirical}

The primary empirical window is the world series in NFA Edition 2026, results and metadata version 1.0. We exclude estimated values for 2024--2025. The 1961--2022 sensitivity uses the same account edition and population series.

From 1961 to 2023, total biocapacity account value $B$ rises from $9.7254\times10^9$ to $12.1428\times10^9$\,gha ($24.86\%$; $\Delta\ln B=+0.221997$). Population rises from $3.06487$ to $8.09173$ billion. Over the same period, $B/P$ falls from $3.17319$ to $1.50065$\,gha\,person$^{-1}$ ($52.71\%$; $\Delta\ln(B/P)=-0.748840$). The Global Ecological Footprint of Production per person, $E_{\mathrm{prod}}/P$, rises from $2.19700$ to $2.56462$\,gha\,person$^{-1}$ ($16.73\%$). The ratio $E_{\mathrm{prod}}/B$ rises from $0.69236$ to $1.70901$ and first exceeds one in 1973.

The exact identity $\Delta\ln B=\Delta\ln P+\Delta\ln(B/P)$ gives $+0.970837-0.748840=+0.221997$. It separates population-scale change from change in biocapacity per person; it does not explain their causes. For 1961--2022, the same-edition sensitivity gives $\Delta\ln B=+0.219497$, a $52.41\%$ fall in $B/P$, a $17.70\%$ rise in $E_{\mathrm{prod}}/P$, and the same year when $E_{\mathrm{prod}}/B$ first exceeds one (SI~\S S5).

\begin{figure}[!htbp]
\centering
\includegraphics[width=\textwidth]{fig2_nfa_aggregation_face}
\caption{Aggregate and per-capita movements in NFA Edition 2026 (world, 1961--2023; complete observations). $B$ is total biocapacity, $E_{\mathrm{prod}}$ is the workbook's Global Ecological Footprint of Production, and $P$ is the single UN PopStat population series; per-capita values are calculated from totals. (a) Indices (1961 $=1$) of $B$, $B/P$, and $E_{\mathrm{prod}}/P$. (b) $E_{\mathrm{prod}}/B$, first above one in 1973. (c) Exact log identity $\Delta\ln B=\Delta\ln P+\Delta\ln(B/P)$. Account values are in gha; plotted indices, ratios, and log changes are dimensionless. The 2024--2025 estimated rows are excluded.}
\label{fig:aggregation}
\end{figure}

The land-type components make the NFA account composition explicit (Fig.~\ref{fig:composition}; Table~\ref{tab:nfa_components}). Cropland biocapacity contributes $+24.06$ percentage points to the $+24.86\%$ aggregate change, about $96.8\%$ of the net percentage change. Forest biocapacity contributes $-5.95$ points, grazing land $-0.27$, fishing grounds $-0.42$, and built-up land $+7.44$. The non-cropland account sum, which includes built-up land, is nearly flat ($1.0092\times$, or $+0.80$ points) only because it nets the built-up increase against declines in other categories. This is measured composition of NFA account values. It does not establish that forest biomass, soil condition, wetland extent, fish biomass, or the model's $A_c$ declined; built-up land is an accounting class, not ecological capital. Component values sum to total biocapacity within workbook rounding at both endpoints.

\begin{table}[!htbp]
\centering\footnotesize
\caption{NFA Edition 2026 biocapacity-account endpoint ratios and contributions. Contributions are $100\,[B_i(T)-B_i(1961)]/B(1961)$ percentage points; endpoint ratios are $B_i(T)/B_i(1961)$. The 2022 sensitivity is from the same edition. All entries are account values, not physical-stock measurements.}
\label{tab:nfa_components}
\begin{tabularx}{\textwidth}{@{}Xrrrr@{}}
\toprule
Component & \shortstack{2023\\ratio} & \shortstack{2023 contribution\\(pp of $B_{1961}$)} & \shortstack{2022 sensitivity\\ratio} & \shortstack{2022 contribution\\(pp of $B_{1961}$)} \\
\midrule
Cropland & 2.8701 & $+24.0577$ & 2.8514 & $+23.8183$ \\
Grazing land & 0.9828 & $-0.2688$ & 0.9815 & $-0.2890$ \\
Forest biocapacity & 0.8996 & $-5.9525$ & 0.9019 & $-5.8137$ \\
Fishing grounds & 0.9627 & $-0.4212$ & 0.9596 & $-0.4559$ \\
Built-up land$^{\dagger}$ & 8.7935 & $+7.4416$ & 8.6298 & $+7.2854$ \\
Non-cropland sum$^{\dagger}$ & 1.0092 & $+0.7990$ & 1.0083 & $+0.7267$ \\
\midrule
Total biocapacity & 1.2486 & $+24.8568$ & 1.2454 & $+24.5450$ \\
\bottomrule
\end{tabularx}
\begin{flushleft}\footnotesize
$^{\dagger}$Built-up land is not ecological capital; the non-cropland account sum includes it. Component sums match total biocapacity within workbook rounding. \end{flushleft}
\end{table}

\begin{figure}[!htbp]
\centering
\includegraphics[width=0.95\textwidth]{fig3_composition_attribution}
\caption{NFA Edition 2026 land-type biocapacity-account changes, world, 1961--2023. (a) Component contributions in percentage points of $B_{1961}$. (b) Endpoint ratios $B_i(2023)/B_i(1961)$. “Forest biocapacity” is the workbook's exact label. All bars are account values, not measurements of physical stocks; built-up land is an NFA accounting class, not ecological capital.}
\label{fig:composition}
\end{figure}

A separate check uses external cropland area and cereal yield to construct a proxy for the fast-land contribution (SI~\S S5). With the 1961 NFA cropland share as the declared anchor, an exact log-change accounting split assigns $+0.268662$ to the cereal-yield index, $+0.037710$ to cropland area, and $-0.084375$ to the residual; the three terms sum to $\Delta\ln B=+0.221997$. This split is exact for the defined proxy, but it changes with the anchor and is not a physical-stock attribution.

\subsection{Falsifiable predictions}\label{sec:predictions}

The model yields the following testable predictions. Tests of land-class composition require observations of aggregate biocapacity, capital-land area, fast-land area, and relevant land-cover or yield proxies, rather than aggregate biocapacity alone.
\begin{enumerate}[label=P\arabic*.,leftmargin=2.2em]
\item \textbf{When the total can rise during conversion.} In the baseline frozen-demand run, a finite interval has $\dot A_c<0$ while $\dot B>0$ (Fig.~\ref{fig:model}a). The conversion contribution is positive exactly when $\dconv>0$, but other land flows, quality, debt, technology, and restoration determine the net change in $B$. With baseline availability shares $\sigma_f=\sigma_c=1$, conversion requires demand to exceed total biocapacity, or $R_B=E/B>1$. With reserved shares below one, the extension can permit conversion even when $R_B\le1$; this is not a baseline result. The composition mask is the pattern $B$ rising while $A_c$ falls; the distinct liquidation mask is $E$ steady while $A_c$ falls. The patterns can overlap, and neither label is a synonym for the other. The sign of $\dconv$ describes the conversion contribution, not necessarily the net change in $B$ when other terms are active.
\item \textbf{How long the pattern lasts.} Its duration depends on the slow return rate $\lambda_1$ (and on $\tconv$, $\mu$, $\eta_f$, and $b_c/b_f$); its size depends on $\dconv$ and the initial shortfall. It is not set by a universal time horizon. In the baseline run, most conversion occurs within one to two e-folding times—the time for the remaining gap to shrink by a factor of $e$.
\item \textbf{No fold-type early warning in the reduced model.} The return time does not grow as demand approaches the ceiling, and the depth of drawdown changes continuously with demand (Fig.~\ref{fig:model}c). Instead, the modeled signal is $A_c$ approaching its minimum. The model does not predict a fold-related warning before this boundary is reached.
\item \textbf{A population threshold with weak delay sensitivity in the tested range.} The initial population at which a run approaches the capital-land minimum rises with starting capital-land area (Table~\ref{tab:model}b). At $A_{c0}=1.0$\,ha, the $10^{-6}$\,ha proximity value changes by at most $0.001$\,cap across $\tau_g=5$, $20$, and $40$\,yr (SI Table S2). This finite-range comparison does not establish delay invariance; the small change is consistent with the modest regenerative contribution and monotone quality trajectory in the baseline.
\item \textbf{Recovery needs a return flow.} Reducing demand alone does not restore capital-land area. An explicit restoration flow is required; a surplus-gated flow is inactive during shortfall, while a deficit-gated flow can deepen the provision deficit.
\end{enumerate}
The predictions cannot be tested from aggregate $B$ alone. Tests of the composition mechanism require independent observations or defensible proxies for $A_c$ and fast-land yield $b_f$; tests of the delay, recovery, and demand-threshold predictions require corresponding trajectories and intervention or demographic data.
\clearpage
\section{Discussion}\label{sec:discussion}

\subsection{What the two-land model shows}

The model yields three results relevant to representations of human--environment systems \citep{Schaefer1954,May1973,Meadows1972,Meadows2004,WackernagelRees1996,BranderTaylor1998}. First, conversion can add to total biocapacity even while capital-land area $A_c$ falls; other land flows and quality change determine whether $B$ rises overall. The exact condition for conversion's contribution is in Proposition~\ref{prop:mask}. The technology boost has a fixed upper limit, while debt depends on the size and duration of shortfall (Proposition~\ref{prop:bounded}).

Second, in the reduced fixed-demand comparison ($q=q_{\max}$, $D=0$, no technology boost, and restoration off), the balance relation decreases steadily with $A_c$ until it meets the minimum-area boundary. It does not turn back to form a fold. This finding—and the absence of hysteresis or critical slowing down—applies to that reduced equilibrium relation, not to every trajectory. With one-way conversion and no return flow, $A_c$ remains below its pre-overshoot level after conversion stops, even though it is no longer declining at the balance point (Table~\ref{tab:neg}). Third, recovery requires an explicit restoration flow, and the effect depends on whether its gate opens during a shortfall or a surplus.

SI~\S S7 presents a one-stock model as a separate comparator. Its state $A$ uses comparator-specific units and is not the capital-land area $A_c$. Without debt and with constant parameters, its equilibrium relation is $B(A)=(b+b_G\rho)A-(b_G\rho/A_{\max})A^2$, with $P=B(A)/e$. At the illustrated comparator settings, $b_G\rho<b$, so the relation is monotone; an interior maximum would require different parameter values. The comparator's neutral equilibrium family and selected delay trajectories apply only to that formulation. They do not establish a fold, a universal critical delay, or a collapse mechanism in the two-land model.

These results narrow the conditions under which common early-warning indicators should be expected. Rising autocorrelation and variance are associated with approach to a fold \citep{Scheffer2001}. In this reduced model, drawdown reaches a minimum-area boundary instead, so those indicators need not rise first. Their absence would not show that the system is safe; the relevant signal is the path of $A_c$, which the aggregate alone cannot reveal.

\begin{table}[!htbp]
\centering\footnotesize
\caption{Equilibrium, trajectory, and tested delay-sensitivity properties under the stated assumptions.}
\label{tab:neg}
\begin{tabularx}{\textwidth}{>{\raggedright\arraybackslash}p{0.34\textwidth} l >{\raggedright\arraybackslash}X}
\toprule
Property & Finding & Scope \\
\midrule
Fold on the reduced frozen-demand branch & Absent & $B_{eq}(A_c)$ is strictly monotone; the balance point collides with the typed floor at $\Eceil\approx1.138\,\mathrm{gha}$ (Proposition~\ref{prop:ceiling}) \\
Critical slowing as $E\to\Eceil$ & Absent & The slow eigenvalue remains negative and constant along the reduced branch \\
Hysteresis of the frozen-demand equilibrium relation & Absent & There is at most one feasible balance point; this does not imply history independence of realized states \\
Path dependence of realized $A_c$ & Present, conditional & With $R_{rc}=0$, one-way conversion leaves $A_c$ below its pre-overshoot level after shortfall closure, while $\dot A_c=0$ at equilibrium \\
Discontinuous response to endogenous demand & Not observed & In the reported runs, minimum $A_c$ varies continuously and monotonically with initial population; proximity thresholds depend on tolerance \\
Regeneration-delay sensitivity ($\tau_g$, with $\tau_p=25$\,yr) & Small in tested range & At $A_{c0}=1.0$\,ha, the delayed-quality $10^{-6}$\,ha threshold varies by $0.001$\,cap across $\tau_g=5,20,40$\,yr (SI Table S2); this does not establish invariance outside that range \\
\bottomrule
\end{tabularx}
\end{table}
\clearpage

\subsection{Measurement implications}

Biocapacity in gha is a value-weighted total, not an inventory of its land classes. A $20\%$ rise in the total does not tell us whether all components grew moderately or whether a large gain in one offset losses in others. Proposition~\ref{prop:mask}(iii) gives the formal limit: at fixed per-hectare yields, one total and the conserved-area constraint still leave many feasible combinations of fast, capital, and reserve land.

Related transition-assessment work defines a linear compensatory operator and tests perfect substitutability against less-substitutable CES alternatives \citep[Sections~2 and 4 of][]{Abaee2026a}. ECOMOD's $B$ is also a linear sum of land contributions. In the companion's separate finite-plan setting, the compensatory protocol accepts states for which no single available plan meets every component floor (Theorem~2 in \citet{Abaee2026a}); no positive-elasticity member is uniformly safe over the specified acceptance-gap region (Theorem~4(iii) in \citet{Abaee2026a}). If feasibility must hold throughout a transition, a temporary fall below a land-class minimum remains a violation even if the endpoint later recovers.

For the National Footprint and Biocapacity Accounts, these findings describe account composition, not physical stocks. In Edition 2026, total world biocapacity account value $B$ rises $24.86\%$ from 1961 to 2023, while $B/P$ falls $52.71\%$. The Global Ecological Footprint of Production per person, $E_{\mathrm{prod}}/P$, rises $16.73\%$, and $E_{\mathrm{prod}}/B$ first exceeds one in 1973. Here $B$ and $E_{\mathrm{prod}}$ are annual account amounts in gha, with the annual basis implicit; the per-person quantities are in gha per person, and $E_{\mathrm{prod}}/B$ is dimensionless.

Cropland contributes $+24.06$ percentage points to the total change; Forest biocapacity contributes $-5.95$, grazing land $-0.27$, fishing grounds $-0.42$, and built-up land $+7.44$. The non-cropland account sum is nearly flat because these changes offset one another, including the large rise in built-up land. The same-edition 1961--2022 comparison is reported in SI~\S S5. These are changes in account values, not measurements of forest biomass or condition, soil carbon, wetland extent, fish biomass, or other physical stocks. In particular, the decline in Forest biocapacity does not by itself establish a decline in forest biomass. Built-up land is an NFA accounting class, not ecological capital. Component ratios, contributions, vintage, and endpoint sensitivity characterize account composition but do not establish physical-stock change \citep{Lo2026NFA,Borucke2013,Blomqvist2013,GiampietroSaltelli2014,Galli2016,vdBerghGrazi2015}.

The model result is separate from the NFA account trends. In the frozen-demand baseline (Fig.~\ref{fig:model}a), $A_c$ falls from $1.500$ to $0.807659$\,ha ($46.16\%$), while model biocapacity $B$ rises from $0.594$ to $0.825$\,gha ($39\%$). Although the quality state $q$ rises from $0.9$ to $1.0$, it starts above $q^*=q_{\max}/2$, so $\Ups(q)$ falls from $0.05576$ to $0.05000$\,gha\,ha$^{-1}$ ($10.33\%$). The modeled contribution $\Ups(q)A_c$ falls from $0.08364$ to $0.04038$\,gha ($51.72\%$).

With restoration disabled, one-way conversion leaves $A_c$ below its starting value after the shortfall closes. The area then persists at its lower balance level; it does not continue to decline indefinitely. These are conditional model outputs, not observations or estimates of global physical-stock change.

We did not assemble a common-boundary global physical-stock series for the relevant components, so we report no physical-stock composite or global decline finding. The SI scoreboard is a specification of candidate stock, extent, condition, and annual-flow measures and the requirements for comparing them; it is not an empirical inventory and contains no endpoint values. Until spatial and temporal boundaries, units, uncertainty, overlap, valuation or normalization, and permitted substitution are specified, these heterogeneous measures should remain a nonadditive vector. SEEA Ecosystem Accounting likewise distinguishes ecosystem extent, condition, services, and monetary presentations \citep{UN2021SEEA_EA}.

The external cropland-area and cereal-yield series is an illustrative proxy, not a direct measurement of the model's land classes. Its 1961 starting value is anchored to the NFA cropland share; cereal yield is a sentinel measure, not an all-crops measure, and the area and account boundaries do not match by construction. Here external cropland area $A_f$ is in ha and the index-scaled yield $\hat b_f$ is in gha\,ha$^{-1}$, so $X=\hat b_fA_f$ and the remainder $C=B-X$ are account amounts in gha. That remainder can include boundary and calibration differences, other account components, and conversion-factor changes; it is neither $\Ups(q)A_c$ nor a physical-stock measure. A physical interpretation would require independent, boundary-consistent data on area, condition, yield, and regeneration, plus a validated mapping and valuation. Forest area may inform one part of the effective capital class but cannot represent the whole class or its condition.

\subsection{Policy controls and limits}

The policy analysis compares three kinds of control. A conversion stop triggered at the minimum capital-land area $A_c=\Acmin$ preserves that floor. A ceiling on fast-land area $A_f$ does not guarantee the capital-land minimum unless the stop rule is explicitly linked to it and to the remaining-area limits. A heuristic governor that monitors only aggregate $B$ cannot guarantee the capital-land floor, because the total does not reveal $A_c$.

The restoration gates also differ: surplus-gated restoration is inactive during shortfall, while a deficit-gated flow acts during shortfall and can reduce provision further. These outcomes connect to the distinction between protective and mobilizing institutional feedback in related work \citep{Abaee2026b}. With a common reserved share $\sigma_f=\sigma_c=\sigma<1$, conversion is inactive while $E\le\sigma B$; with distinct shares, the condition is $E\le\sigma_fY_f+\sigma_cY_c$. This share rule limits the human-available fraction of annual account supply; it is not a spatial map of protected land or an evaluation of a specific conservation proposal. Robustness to disturbances and parameter uncertainty has not been established.

\subsection{Relation to prior work}

The ledger companion develops component-level accounting for material flows, and the assessment companion examines how aggregate measures allow one component to compensate for another \citep{Abaee2026a,Abaee2026c}. This paper adds a delayed land-class model for biocapacity in global hectares. The governance companion studies institutional feedback \citep{Abaee2026b}; the delays here represent ecological regeneration and demographic response, not institutional review.

\subsection{Limitations}

Several limitations govern how to read the model. Ecological-capital land combines diverse assets, so its single quality state, regeneration rule, composition diagnostic, and demand ceiling are effective model quantities—not direct estimates for forests, soils, or wetlands. Most parameters are representative rather than estimated. A cross-section cannot separately identify $b_c$, $b_{G,c}$, and $\rho_c$; recovery records from different systems provide timescale context, not a global calibration of $\rho_c$. They describe regeneration response, not the delay $\tau_g$. The per-person footprint is held constant to isolate the feedbacks among land, biocapacity, and demand; allowing it to change is a further extension.

The composition pattern depends on demand and the mechanisms specified here; under these assumptions, it is bounded. The exact shortfall switch makes stability at $S=0$ one-sided. Reported drawdown thresholds are proximity bands because their values depend on the chosen tolerance and time step. The debt-repayment rate is fixed at $\eta=0.05$\,yr$^{-1}$; the zero-repayment limit, $\eta\to0$, removes the debt equilibrium and is not analyzed. NFA accounts combine area and yield, so they cannot identify those quantities separately. Finally, using current-quality regeneration $G_c(q(t))$ rather than the delayed form $G_c(q(t-\tau_g))$ lengthens the demonstration mask by $1.3$\,yr and shifts the initial-population threshold by at most $0.001$\,cap (SI~\S S3).

\section{Conclusions}\label{sec:conclusions}

Aggregate biocapacity does not reveal how its land classes are changing. In the two-land model, conversion adds to $B$ when fast-land yield exceeds the effective per-hectare value of capital land, even as $A_c$ falls; other terms determine the net change in $B$ (Proposition~\ref{prop:mask}). Under the reduced fixed-demand assumptions, the balance relation has one branch and reaches a minimum-area boundary without a fold, hysteresis, or critical slowing down. The realized path is still history-dependent: with restoration off, $A_c$ remains below its pre-overshoot level after the shortfall closes, but does not continue declining at the balance point. In the representative run, $A_c$ falls $46.16\%$, $\Ups(q)$ falls $10.33\%$, and $\Ups(q)A_c$ falls $51.72\%$ while model $B$ rises $39\%$. These are conditional model results, not estimates of global physical stocks.

In NFA Edition 2026, total world biocapacity account value rises $24.86\%$ from 1961 to 2023, while $B/P$ falls $52.71\%$ and the production footprint per person rises $16.73\%$. Cropland adds $24.06$ percentage points to the account total; Forest biocapacity subtracts $5.95$ points, and built-up land adds $7.44$ as an accounting class, not ecological capital. These account movements are not measurements of biomass, soil condition, wetland extent, fish biomass, or $A_c$. The 1961--2022 sensitivity uses the same account edition and population series. No common-boundary global physical-stock series for the relevant components was assembled, so we report no physical-stock composite or global decline finding. The SI scoreboard is a specification, not an empirical inventory; heterogeneous physical measures remain a nonadditive vector until boundaries, units, uncertainty, overlap, valuation or normalization, and allowed substitution are made explicit. The proxy residual $C=B-\hat b_fA_f$ is neither a physical-capital measure nor the model contribution $\Ups(q)A_c$. A future physical-stock comparison would require independent observations and a validated mapping.

The practical reporting lesson is to publish component accounts alongside the aggregate. The orchard is a metaphor for this point, not evidence of a measured global stock decline. In the analogy, one book records the annual harvest and a second records the trees, soil, and other assets that support later harvests. A real physical-stock claim would require that second book to be filled with independent, comparable observations—not inferred from the account total. A harvest total is not a census of the orchard.

\section*{Acknowledgments}

The author thanks Global Footprint Network for making its data publicly available; Mathis Wackernagel for his thoughtful review; the editors of \emph{Ecological Modelling} for their trust and support; and Arena.ai for its ingenuity and online tools.

\section*{Declarations}

\paragraph{Data availability} All empirical inputs are publicly available from the sources cited in the manuscript. Derived analysis data and model code will be deposited in a public GitHub repository and archived in Zenodo; repository and record details will be added after the deposit is complete.

\paragraph{Funding} No funding was received.

\paragraph{Competing interests} The author declares no competing interests.

\paragraph{CRediT authorship contribution statement} \textbf{Amin Abaee:} Conceptualization; Methodology; Software; Formal analysis; Writing -- original draft; Writing -- review \& editing.

\paragraph{Declaration of generative AI in scientific writing} DeepSeek AI and Claude (Anthropic) assisted with computational-reproducibility checks under the author's supervision; the author reviewed and edited the content and takes full responsibility for the publication.
\begin{thebibliography}{99}
\bibitem[{Abaee(2026a)}]{Abaee2026a} Abaee, A. (2026a). \emph{Aggregation is a claim, not a presentation: the geometry of the acceptance gap between compensatory and noncompensatory transition assessment}. Working paper (v71). \url{https://github.com/MIKEAA2020/general-sustainability/blob/2efaa62083f2007dec3446c15c475af88f0925f3/arena%20agent%201/paper%20rewrites/latex/paper06_assessment_separation_v71.pdf}.
\bibitem[{Abaee(2026b)}]{Abaee2026b} Abaee, A. (2026b). \emph{Governance delay and the stability of harvested stocks: mobilising and protective feedback rules, and the review interval as a design parameter}. Working paper (v41). \url{https://github.com/MIKEAA2020/general-sustainability/blob/2efaa62083f2007dec3446c15c475af88f0925f3/arena%20agent%201/paper%20rewrites/latex/paper4_delay_dynamics_v41.pdf}.
\bibitem[{Abaee(2026c)}]{Abaee2026c} Abaee, A. (2026c). \emph{Typed Flux Ledgers and Depletion Arithmetic: Conservation, Componentwise Diagnostics, and the Semantics of Depletion Horizons}. Working paper (v50). \url{https://github.com/MIKEAA2020/general-sustainability/blob/2efaa62083f2007dec3446c15c475af88f0925f3/paper3/latest-2026-09-19/manuscript/paper3_material_ledgers_v50.pdf}.
\bibitem[{Anderies(2000)}]{Anderies2000} Anderies, J. M. (2000). On modeling human behavior and institutions in simple ecological economic systems. \emph{Ecological Economics}, 35(3), 393--412.
\bibitem[{Becker et al.(2017)}]{Becker2017} Becker, W., Saisana, M., Paruolo, P. \& Vandecasteele, I. (2017). Weights and importance in composite indicators: Closing the gap. \emph{Ecological Indicators}, 80, 12--22.
\bibitem[{Blomqvist et al.(2013)}]{Blomqvist2013} Blomqvist, L., Brook, B. W., Ellis, E. C., Kareiva, P. M., Nordhaus, T. \& Shellenberger, M. (2013). Does the shoe fit? Real versus imagined ecological footprints. \emph{PLoS Biology}, 11(11), e1001700.
\bibitem[{Borucke et al.(2013)}]{Borucke2013} Borucke, M., Moore, D., Cranston, G., Gracey, K., Iha, K., Larson, J., Lazarus, E., Morales, J. C., Wackernagel, M. \& Galli, A. (2013). Accounting for demand and supply of the biosphere's regenerative capacity: The National Footprint Accounts' underlying methodology and framework. \emph{Ecological Indicators}, 24, 518--533.
\bibitem[{Brander and Taylor(1998)}]{BranderTaylor1998} Brander, J. A. \& Taylor, M. S. (1998). The simple economics of Easter Island: A Ricardo--Malthus model of renewable resource use. \emph{American Economic Review}, 88(1), 119--138.
\bibitem[{Calvin and Bond-Lamberty(2018)}]{CalvinBondLamberty2018} Calvin, K. \& Bond-Lamberty, B. (2018). Integrated human--Earth system modeling --- state of the science and future directions. \emph{Environmental Research Letters}, 13(6), 063006.
\bibitem[{Cohen(1995)}]{Cohen1995} Cohen, J. E. (1995). Population growth and Earth's human carrying capacity. \emph{Science}, 269(5222), 341--346.
\bibitem[{Dasgupta(2021)}]{Dasgupta2021} Dasgupta, P. (2021). \emph{The Economics of Biodiversity: The Dasgupta Review.} London: HM Treasury.
\bibitem[{Donges et al.(2017)}]{Donges2017} Donges, J. F., Winkelmann, R., Cornell, S. E., Lucht, W., Dyke, J. G., Rockstr\"om, J., Heitzig, J. \& Schellnhuber, H. J. (2017). Closing the loop: Reconnecting human dynamics to Earth system science. \emph{The Anthropocene Review}, 4(2), 151--157.
\bibitem[{Erb et al.(2017)}]{Erb2017} Erb, K.-H., Luyssaert, S., Meyfroidt, P., Pongratz, J., Don, A., Kloster, S., Kuemmerle, T., Fetzel, T., Fuchs, R., Hasegawa, T., Haberl, H., et al. (2017). Land management: data availability and process understanding for global change studies. \emph{Global Change Biology}, 23(2), 512--533.
\bibitem[{Fischer et al.(2022)}]{Fischer2022} Fischer, S. M., Joy, M. K., Abrahamse, W., Milfont, T. L. \& Petherick, L. M. (2022). The use and misuse of composite environmental indices. \emph{bioRxiv}, 2022.03.15.484501. \url{https://doi.org/10.1101/2022.03.15.484501}.
\bibitem[{Global Footprint Network(2009)}]{GFNStandards2009} Global Footprint Network (2009). \emph{Ecological Footprint Standards 2009}. Global Footprint Network. \url{https://www.footprintnetwork.org/content/uploads/2019/05/Ecological_Footprint_Standards_2009.pdf}.
\bibitem[{Global Footprint Network(2019)}]{GFNGuidebook2019} Global Footprint Network (2019). \emph{Working Guidebook to the National Footprint Accounts 2019}. Global Footprint Network. \url{https://www.footprintnetwork.org/content/uploads/2019/05/National_Footprint_Accounts_Guidebook_2019.pdf}.
\bibitem[{Galli et al.(2016)}]{Galli2016} Galli, A., Giampietro, M., Goldfinger, S., Lazarus, E., Lin, D., Saltelli, A., Wackernagel, M. \& M\"uller, F. (2016). Questioning the ecological footprint. \emph{Ecological Indicators}, 69, 224--232.
\bibitem[{Giampietro and Saltelli(2014)}]{GiampietroSaltelli2014} Giampietro, M. \& Saltelli, A. (2014). Footprints to nowhere. \emph{Ecological Indicators}, 46, 610--621.
\bibitem[{Haberl et al.(2007)}]{Haberl2007} Haberl, H., Erb, K.-H., Krausmann, F., Gaube, V., Bondeau, A., Plutzar, C., Gingrich, S., Lucht, W. \& Fischer-Kowalski, M. (2007). Quantifying and mapping the human appropriation of net primary production in earth's terrestrial ecosystems. \emph{PNAS}, 104(31), 12942--12947.
\bibitem[{Henderson and Loreau(2018)}]{HendersonLoreau2018} Henderson, K. \& Loreau, M. (2018). How ecological feedbacks between human population and land cover influence sustainability. \emph{PLoS Computational Biology}, 14(8), e1006389.
\bibitem[{Hutchings and Reynolds(2004)}]{HutchingsReynolds2004} Hutchings, J. A. \& Reynolds, J. D. (2004). Marine fish population collapses: consequences for recovery and extinction risk. \emph{BioScience}, 54(4), 297--309.
\bibitem[{Hutchinson(1948)}]{Hutchinson1948} Hutchinson, G. E. (1948). Circular causal systems in ecology. \emph{Annals of the New York Academy of Sciences}, 50(4), 221--246.
\bibitem[{Krausmann et al.(2013)}]{Krausmann2013} Krausmann, F., Erb, K.-H., Gingrich, S., Haberl, H., Bondeau, A., Gaube, V., Lauk, C., Plutzar, C. \& Searchinger, T. D. (2013). Global human appropriation of net primary production doubled in the 20th century. \emph{PNAS}, 110(25), 10324--10329.
\bibitem[{Lo et al.(2026)}]{Lo2026NFA} Lo, K., Basnet, N., Miller, E., Srikanthalingam, B., Foley, B., Van Berkum, J. L., Long, A. H., Toneva, P., Ermina, M. \& Klinkenberg, C. (2026). \emph{National Ecological Footprint and Biocapacity Accounts, Edition 2026 (1961--2025): Results and metadata, version 1.0.} Prepared for Footprint Data Foundation by researchers at York University and the University of Iceland. Workbook: \url{https://footprint.info.yorku.ca/files/2026/04/NEFBA_Data_global_national_ecological_footprint_biocapacity.xlsx}; release notes: \url{https://footprint.info.yorku.ca/files/2026/04/NEFBA_ReleaseNotes.pdf}; data page: \url{https://footprint.info.yorku.ca/data/}.
\bibitem[{May(1973)}]{May1973} May, R. M. (1973). \emph{Stability and Complexity in Model Ecosystems.} Princeton: Princeton University Press.
\bibitem[{Meadows et al.(1972)}]{Meadows1972} Meadows, D. H., Meadows, D. L., Randers, J. \& Behrens, W. W. (1972). \emph{The Limits to Growth.} New York: Universe Books.
\bibitem[{Meadows et al.(1974)}]{Meadows1974} Meadows, D. L., Behrens, W. W., Meadows, D. H., Naill, R. F., Randers, J. \& Zahn, E. (1974). \emph{Dynamics of Growth in a Finite World.} Cambridge, MA: Wright-Allen Press.
\bibitem[{Meadows et al.(2004)}]{Meadows2004} Meadows, D. H., Randers, J. \& Meadows, D. L. (2004). \emph{The Limits to Growth: The 30-Year Update.} White River Junction, VT: Chelsea Green.
\bibitem[{Meyer(1994)}]{Meyer1994} Meyer, P. S. (1994). Bi-logistic growth. \emph{Technological Forecasting and Social Change}, 47(1), 89--102.
\bibitem[{Meyer and Ausubel(1999)}]{MeyerAusubel1999} Meyer, P. S. \& Ausubel, J. H. (1999). Carrying capacity: A model with logistically varying limits. \emph{Technological Forecasting and Social Change}, 61(3), 209--214.
\bibitem[{Motesharrei et al.(2016)}]{Motesharrei2016} Motesharrei, S., Rivas, J., Kalnay, E., Asrar, G. R., Busalacchi, A. J., Cahalan, R. F., Cane, M. A., Colwell, R. R., Feng, K., Franklin, R. S., et al. (2016). Modeling sustainability: Population, inequality, consumption, and bidirectional coupling of the Earth and human systems. \emph{National Science Review}, 3(4), 470--494.
\bibitem[{Nardo et al.(2008)}]{Nardo2008} Nardo, M., Saisana, M., Saltelli, A., Tarantola, S., Hoffmann, A. \& Giovannini, E. (2008). \emph{Handbook on Constructing Composite Indicators: Methodology and User Guide.} Paris: OECD Publishing.
\bibitem[{Neubauer et al.(2013)}]{Neubauer2013} Neubauer, P., Jensen, O. P., Hutchings, J. A. \& Baum, J. K. (2013). Resilience and recovery of overexploited marine populations. \emph{Science}, 340(6130), 347--349.
\bibitem[{Our~World~in~Data(2026a)}]{OWIDCereal2026} Our~World~in~Data (2026a). \emph{Cereal yield.} FAOSTAT data accessed through Our World in Data; grapher page last updated 25 February 2026; accessed 24 October 2026. \url{https://ourworldindata.org/grapher/cereal-yield}.
\bibitem[{Our~World~in~Data(2026b)}]{OWIDCropland2026} Our~World~in~Data (2026b). \emph{Cropland area.} FAOSTAT data accessed through Our World in Data; grapher page last updated 25 February 2026; accessed 24 October 2026. \url{https://ourworldindata.org/grapher/cropland-area}.
\bibitem[{Pearl(1925)}]{Pearl1925} Pearl, R. (1925). \emph{The Biology of Population Growth.} New York: Alfred A. Knopf.
\bibitem[{Poeplau et al.(2011)}]{Poeplau2011} Poeplau, C., et al. (2011). Temporal dynamics of soil organic carbon after land-use change in the temperate zone--carbon response functions as a model approach. \emph{Global Change Biology}, 17(7), 2415--2427.
\bibitem[{Poorter et al.(2016)}]{Poorter2016} Poorter, L., et al. (2016). Biomass resilience of Neotropical secondary forests. \emph{Nature}, 530(7590), 211--214.
\bibitem[{Schaefer(1954)}]{Schaefer1954} Schaefer, M. B. (1954). Some aspects of the dynamics of populations important to the management of the commercial marine fisheries. \emph{Bulletin of the Inter-American Tropical Tuna Commission}, 1(2), 25--56.
\bibitem[{Scheffer et al.(2001)}]{Scheffer2001} Scheffer, M., Carpenter, S., Foley, J. A., Folke, C. \& Walker, B. (2001). Catastrophic shifts in ecosystems. \emph{Nature}, 413(6856), 591--596.
\bibitem[{Simon(1996)}]{Simon1996} Simon, J. L. (1996). \emph{The Ultimate Resource 2.} Princeton: Princeton University Press.
\bibitem[{United Nations et al.(2021)}]{UN2021SEEA_EA} United Nations et al. (2021). \emph{System of Environmental-Economic Accounting---Ecosystem Accounting (SEEA EA).} White cover publication, pre-edited text subject to official editing. New York: United Nations. \url{https://seea.un.org/ecosystem-accounting}.
\bibitem[{United Nations(2022)}]{UN2022} United Nations (2022). \emph{World Population Prospects 2022: Summary of Results.} New York: Department of Economic and Social Affairs, Population Division.
\bibitem[{van den Bergh and Grazi(2015)}]{vdBerghGrazi2015} van den Bergh, J. C. J. M. \& Grazi, F. (2015). Reply to the first systematic response by the Global Footprint Network to criticism: a real debate finally? \emph{Ecological Indicators}, 58, 458--463.
\bibitem[{van den Bergh and Rietveld(2004)}]{vdBerghRietveld2004} van den Bergh, J. C. J. M. \& Rietveld, P. (2004). Reconsidering the limits to world population: Meta-analysis and meta-prediction. \emph{BioScience}, 54(3), 195--204.
\bibitem[{Wackernagel and Rees(1996)}]{WackernagelRees1996} Wackernagel, M. \& Rees, W. (1996). \emph{Our Ecological Footprint: Reducing Human Impact on the Earth.} New Society Publishers.
\end{thebibliography}

\end{document}
